\documentclass[10pt,a4paper,reqno]{amsart}

\usepackage[T1]{fontenc}
\usepackage[utf8]{inputenc}
\usepackage{newtxtext,newtxmath}
\usepackage[final]{microtype}
\usepackage{mathtools}
\usepackage{amscd}
\usepackage{xcolor}
\usepackage{hyperref}

\providecommand{\paperauthor}{FDmd-233}

\hypersetup{
  unicode=true,
  pdfencoding=auto,
  colorlinks=true,
  linkcolor=blue!45!black,
  citecolor=blue!45!black,
  urlcolor=blue!45!black,
  pdftitle={Prill-Exceptional Étale Covers from Raynaud Bundles},
  pdfsubject={Finite étale covers of algebraic curves and Raynaud bundles},
  pdfkeywords={Prill's problem, étale covers, Raynaud bundles, finite monodromy},
  pdfauthor={\paperauthor}
}

\newtheorem{theorem}{Theorem}[section]
\newtheorem{lemma}[theorem]{Lemma}
\newtheorem{proposition}[theorem]{Proposition}
\theoremstyle{remark}
\newtheorem{remark}[theorem]{Remark}

\newcommand{\CC}{\mathbb{C}}
\newcommand{\OO}{\mathcal{O}}
\newcommand{\cHom}{\mathcal{H}\!om}
\newcommand{\tens}{\otimes}
\DeclareMathOperator{\Pic}{Pic}
\DeclareMathOperator{\Jac}{Jac}
\DeclareMathOperator{\rk}{rk}
\DeclareMathOperator{\Spec}{Spec}

\title[Prill-exceptional covers]{Prill-Exceptional \'{E}tale Covers from Raynaud Bundles}
\ifx\paperauthor\empty\else
  \author{\paperauthor}
\fi
\date{10 July 2026}

\subjclass[2020]{Primary 14H30; Secondary 14H60, 14H40}
\keywords{Prill's problem, finite \'{e}tale covers, Raynaud bundles, finite monodromy}

\begin{document}

\begin{abstract}
Let $C$ be a smooth projective connected complex curve of genus $g\geq 2$.
We prove that $C$ admits a connected finite \'{e}tale Galois cover
$f\colon X\to C$ such that
\[
  h^0\!\left(X,\OO_X\bigl(f^{-1}(p)\bigr)\right)\geq 2
\]
for every $p\in C$.
Let $V=R_g\tens M$ be a determinant-normalised twist of the rank-$g^g$
Raynaud bundle.  Then $\deg V=0$, $\det V\cong\OO_C$, and $H^0(C,V)=0$,
whereas $H^0(C,V(p))\neq0$ for every $p\in C$.  We prove that $V$ has finite
monodromy and choose an irreducible summand $W$ with
$H^0(C,W(p))\neq0$ for all $p$.  For the Galois cover defined by the
monodromy of $W$, the associated bundle $f_*\OO_X$ contains $\OO_C$ and $W$
as direct summands, which yields the claimed inequality.
\end{abstract}

\maketitle

\section{Introduction}

Let $f\colon X\to C$ be a finite morphism of smooth projective connected
complex curves.  Prill asked whether a fibre divisor $f^{-1}(p)$ can move in
a pencil; equivalently, whether
\[
  h^0\!\left(X,\OO_X\bigl(f^{-1}(p)\bigr)\right)>1
\]
for a general point $p\in C$.  The question appears in
\cite[Chapter~VI, Exercise~D]{ACGH} and, for curves of genus at least three,
is recorded as \cite[Problem~14]{LittProblem14}.  Partial affirmative results
were obtained in \cite{BiswasButler,MohanKumar}.  Landesman and Litt
subsequently constructed, for every curve of genus two, a connected finite
\'{e}tale cover of degree $36$ satisfying the inequality for every $p\in C$
\cite{LandesmanLitt}.

We call $f$ \emph{Prill-exceptional} when the inequality holds for every
$p\in C$.  The construction below works for every $g\geq2$.

\begin{theorem}\label{thm:main}
Let $C$ be a smooth projective connected complex curve of genus $g\geq 2$.
There is a connected finite \'{e}tale Galois cover $f\colon X\to C$ such that
\[
  h^0\!\left(X,\OO_X\bigl(f^{-1}(p)\bigr)\right)\geq 2
\]
for every $p\in C$.
\end{theorem}

For $g\geq 3$, Theorem~\ref{thm:main} answers
\cite[Problem~14]{LittProblem14} affirmatively.  For a finite flat morphism,
the function
\[
  p\longmapsto h^0\!\left(X,\OO_X\bigl(f^{-1}(p)\bigr)\right)
\]
is upper semicontinuous.  The locus where this dimension is at least two is
therefore closed.  If it contains a dense open subset of $C$, it equals $C$.
Thus the general-fibre and pointwise formulations are equivalent.  We use the
latter because the construction yields it directly.

The proof uses two properties of Raynaud's construction.  The restricted
bundle $R_g$ satisfies $H^0(C,R_g\tens\alpha)\neq0$ for every
$\alpha\in\Pic^0(C)$, whereas the ambient bundle $E_g$ becomes a direct sum
of copies of one line bundle after pullback by $[g]\colon J\to J$.  We choose
a twist with trivial determinant while preserving the pointwise
nonvanishing.  On the induced \'{e}tale pullback of $C$, the repeated line
bundle is torsion, so the twisted bundle has finite linear monodromy.

\section{The Raynaud bundle}

Write $J=\Jac(C)=\Pic^0(C)$, fix a theta divisor $\Theta$ on $J$, and choose
an Abel--Jacobi embedding $i\colon C\hookrightarrow J$.  The form of Raynaud's
construction needed here is recorded in \cite[Section~3]{Raynaud} and
\cite[Section~2.3]{Beauville}.

\begin{theorem}[Raynaud]\label{thm:raynaud}
For every integer $n\geq 1$, there is a vector bundle $E_n$ on $J$ such that
\begin{equation}\label{eq:raynaud-pullback}
  [n]^*E_n
  \cong H^0\!\left(J,\OO_J(n\Theta)\right)^*\tens_{\CC}\OO_J(n\Theta).
\end{equation}
Its restriction $R_n=i^*E_n$ is semistable, has rank $n^g$ and slope $g/n$,
and satisfies
\begin{equation}\label{eq:raynaud-nonvanishing}
  H^0(C,R_n\tens\alpha)\neq 0
  \qquad\text{for every }\alpha\in\Pic^0(C).
\end{equation}
\end{theorem}

Take $n=g$ and set
\[
  R:=R_g,
  \qquad
  r:=\rk(R)=g^g.
\]
Then $R$ is semistable of slope one, so $\deg(R)=r$.

\section{A determinant-normalised twist}

\begin{lemma}\label{lem:twist}
There is a line bundle $M\in\Pic^{-1}(C)$ such that
\begin{equation}\label{eq:choice-M}
  M^{\tens r}\cong(\det R)^{-1}
  \qquad\text{and}\qquad
  H^0(C,R\tens M)=0.
\end{equation}
\end{lemma}

\begin{proof}
The map
\[
  \Pic^{-1}(C)\longrightarrow\Pic^{-r}(C),
  \qquad
  M\longmapsto M^{\tens r},
\]
is a translate of the multiplication-by-$r$ isogeny of $J$.  Consequently,
the fibre
\[
  T:=\left\{M\in\Pic^{-1}(C):M^{\tens r}\cong(\det R)^{-1}\right\}
\]
has cardinality $r^{2g}$.

Call $M\in\Pic^{-1}(C)$ \emph{bad} if $H^0(C,R\tens M)\neq 0$.  Such a
section induces a nonzero morphism $M^{-1}\to R$.  If $L\subset R$ is the
saturation of its image, then $L\cong M^{-1}(D)$ for an effective divisor
$D$.  Since $\deg(M^{-1})=1$ and $R$ is semistable of slope one,
\[
  1+\deg D=\deg L\leq 1.
\]
Thus $D=0$, and $M^{-1}$ is a degree-one line subbundle of $R$.

Fix a Jordan--H\"older filtration
$0=R_0\subset R_1\subset\cdots\subset R_m=R$.  Given a degree-one line
subbundle $L\subset R$, let $k$ be minimal with $L\subset R_k$.  The induced
map $L\to R_k/R_{k-1}$ is nonzero.  Since its source and target are stable of
slope one, it is an isomorphism.  Thus every bad $M$ makes $M^{-1}$
isomorphic to a rank-one Jordan--H\"older factor of $R$.  The filtration has
at most $r$ factors, so at most $r$ isomorphism classes of $M$ are bad.

Because $g\geq 2$ and $r=g^g>1$, one has $r^{2g}>r$.  The set $T$ therefore
contains an $M$ that is not bad.
\end{proof}

Fix $M$ as in Lemma~\ref{lem:twist} and put $V:=R\tens M$.  Then
\begin{equation}\label{eq:V-properties}
  \deg V=0,
  \qquad
  \det V\cong\OO_C,
  \qquad
  H^0(C,V)=0.
\end{equation}
For every $p\in C$, the line bundle $M(p)$ has degree zero, so
\eqref{eq:raynaud-nonvanishing} gives
\begin{equation}\label{eq:V-pointwise}
  H^0(C,V(p))=H^0(C,R\tens M(p))\neq 0.
\end{equation}

\section{Finite monodromy}

\begin{proposition}\label{prop:finite-monodromy}
The vector bundle $V$ is trivialised by a connected finite \'{e}tale cover of
$C$.
\end{proposition}

\begin{proof}
Form the cartesian square
\[
\begin{CD}
  \widetilde C @>{\widetilde i}>> J\\
  @V{q}VV                         @VV{[g]}V\\
  C              @>{i}>>          J.
\end{CD}
\]
The morphism $q$ is finite \'{e}tale and surjective.  By
\eqref{eq:raynaud-pullback},
\[
  q^*V
  \cong
  H^0\!\left(J,\OO_J(g\Theta)\right)^*
  \tens_{\CC}
  \bigl(\widetilde i^*\OO_J(g\Theta)\tens q^*M\bigr).
\]

Let $\widetilde C_0$ be a connected component of $\widetilde C$.  Its image
under $q$ is nonempty, open, and closed, hence equals $C$.  Define
\[
  N:=\left.
  \bigl(\widetilde i^*\OO_J(g\Theta)\tens q^*M\bigr)
  \right|_{\widetilde C_0}.
\]
Since $h^0(J,\OO_J(g\Theta))=g^g=r$, the restriction of $q^*V$ to
$\widetilde C_0$ is $N^{\oplus r}$.  Taking determinants and using
$\det V\cong\OO_C$ yields
\[
  N^{\tens r}\cong\OO_{\widetilde C_0}.
\]

Let $d$ be the order of $N$; then $d\mid r$.  A chosen isomorphism
$N^{\tens d}\cong\OO_{\widetilde C_0}$ defines the finite \'{e}tale
$\mu_d$-torsor
\[
  \Spec_{\widetilde C_0}
  \left(\bigoplus_{a=0}^{d-1}N^{-a}\right)
  \longrightarrow \widetilde C_0 .
\]
The pullback of $N$ to this torsor is
trivial.  Any connected component maps surjectively to $\widetilde C_0$; its
composite with $\widetilde C_0\to C$ is a connected finite \'{e}tale cover
$h\colon C'\to C$ that trivialises $V$.
\end{proof}

Passing to a connected finite \'{e}tale Galois cover dominating $h$, we may
regard $V$ as the bundle associated with a complex representation of a finite
group.

\begin{proposition}\label{prop:irreducible-summand}
There is a nontrivial irreducible finite-monodromy vector bundle $W$ on $C$
such that
\[
  H^0(C,W(p))\neq 0
  \qquad\text{for every }p\in C.
\]
\end{proposition}

\begin{proof}
Choose a connected finite \'{e}tale Galois cover $u\colon Y\to C$ trivialising
$V$, and write $H$ for its Galois group.  Fix a trivialisation
$u^*V\cong\OO_Y^{\oplus r}$.  Since $Y$ is connected and projective,
$H^0(Y,\OO_Y)=\CC$; hence the descent matrices are constant and define a
representation $H\to\operatorname{GL}_r(\CC)$.  Maschke's theorem gives a
decomposition
\[
  V\cong\bigoplus_{j=1}^{s}W_j^{\oplus m_j},
\]
where the $W_j$ are pairwise nonisomorphic irreducible finite-monodromy
bundles.  No $W_j$ is trivial, because a trivial summand would contradict
$H^0(C,V)=0$ in \eqref{eq:V-properties}.

For $1\leq j\leq s$, set
\[
  Z_j:=\{p\in C\mid H^0(C,W_j(p))\neq 0\}.
\]
Let $p_1,p_2\colon C\times C\to C$ be the projections and let
$\Delta\subset C\times C$ be the diagonal.  The restriction of
$p_1^*W_j\tens\OO_{C\times C}(\Delta)$ to $C\times\{p\}$ is $W_j(p)$.
Upper semicontinuity for the proper morphism $p_2$ shows that each $Z_j$ is
closed.  The decomposition of $V$ and \eqref{eq:V-pointwise} give
$C=\bigcup_{j=1}^{s}Z_j$.  An irreducible variety cannot be a finite union of
proper closed subsets, so $Z_j=C$ for some $j$.
\end{proof}

\section{The covering}

The vector-bundle criterion used below is the \'{e}tale case of
\cite[Lemma~5.4]{LandesmanLittPW}; we include the argument for completeness.

\begin{proof}[Proof of Theorem~\ref{thm:main}]
Let $W$ be as in Proposition~\ref{prop:irreducible-summand}, and let
\[
  \rho\colon\pi_1(C,c)\longrightarrow\operatorname{GL}(W_c)
\]
be its monodromy representation.  Its image $G$ is finite.  The subgroup
$\ker\rho$ determines a connected finite \'{e}tale Galois cover
$f\colon X\to C$ with Galois group $G$.

The bundle $f_*\OO_X$ corresponds to the regular $G$-module $\CC[G]$, which
contains every irreducible complex representation of $G$ with multiplicity
equal to its dimension.  Hence $\OO_C$ and $W$ are distinct direct summands
of $f_*\OO_X$.

For $p\in C$, the fibre divisor is the pullback of $p$, so
\[
  \OO_X\bigl(f^{-1}(p)\bigr)\cong f^*\OO_C(p).
\]
The projection formula and the summands $\OO_C,W\subset f_*\OO_X$ give
\begin{align*}
  H^0\!\left(X,\OO_X\bigl(f^{-1}(p)\bigr)\right)
  &\cong H^0\!\left(C,f_*\OO_X\tens\OO_C(p)\right)\\
  &\supset H^0(C,\OO_C(p))\oplus H^0(C,W(p)).
\end{align*}
Since $p$ is effective, $h^0(C,\OO_C(p))\geq1$; Proposition~\ref{prop:irreducible-summand}
gives $h^0(C,W(p))\geq1$.  Thus the left-hand side has dimension at least two
for every $p\in C$.
\end{proof}

\begin{remark}
The construction gives no useful degree bound.  The initial Raynaud bundle
has rank $g^g$, and the proof subsequently takes a Kummer cover and a Galois
closure.  In genus two, the degree-$36$ cover of \cite{LandesmanLitt} is
numerically far more efficient.
\end{remark}

\begin{remark}
The determinant condition is essential.  Formula
\eqref{eq:raynaud-pullback} shows only that the pullback of $R_g$ is a direct
sum of copies of one line bundle.  Triviality of the determinant makes that
line bundle torsion.  Consequently, the twisted bundle has finite linear,
rather than merely finite projective, monodromy.
\end{remark}

\bibliographystyle{amsplain}
\bibliography{source/references}

\end{document}
